\subsection{Ablation Study}
\label{sec:ablation}

Matchmaker builds its test inputs from three components. \emph{Boundary-positive
generation} produces strings that the regex accepts, chosen to cover the minimum
and maximum repeat counts, the first and last characters of each class, and every
alternation branch (\S\ref{sec:positive}). Negative generation then runs in two
steps. \emph{Local complements} (step~1) take one representative accepted string
and, at each literal, character class, or backreference position, replace the
character with one that the regex rejects at that position, giving one single-fault
string per position. \emph{Context-aware mutations} (step~2) apply violation
families selected from the regex's structure (wrong-character substitution,
character-class, separator, numeric, quantifier, group, and escape violations,
among others) to the same representative string, cycling through the families
round-robin. If steps~1 and~2 together yield fewer than $N$ distinct rejected
strings, a fallback fill tops the set up to $N$. To measure what each component
contributes, we remove or replace one component at a time and compare the
coverage of the remaining configuration with that of the full system.

\paragraph{Variants.}
We compare the full system with four variants (Table~\ref{tab:ablation}).
\begin{itemize}
  \item \textbf{Full}: boundary positives, followed by steps~1 and~2 and the
    fallback fill for negatives.
  \item \textbf{Shortest-first enumeration}: the boundary positives are replaced by
    the first $N$ distinct accepted strings in length order (shortest first; within
    one length, alternation branches from left to right and class characters in
    code-point order). This baseline is systematic but takes no account of
    boundaries. Negatives are the same as in Full.
  \item \textbf{No positives}: positives are removed, and negatives are the same as
    in Full. This variant shows how much of the coverage the negatives reach alone.
  \item \textbf{Local complements only}: boundary positives, with local complements as the only
    negatives. There are no mutations and no fallback fill.
  \item \textbf{Mutations only}: boundary positives, with context-aware mutations as the
    only negatives. There is no fallback fill, and any string that step~1 would
    produce is removed.
\end{itemize}
Neither step variant includes the fallback fill. The fill always tops the
negative set up to $N$ strings, so it would hide part of the step that was
removed. Mutations only removes step~1's strings because several mutation families
(wrong-character substitution, character-class violation, and case sensitivity
among them) produce exactly the one-character replacements that step~1 makes.
Without this exclusion, Mutations only would still contain most of the local
complements, and the comparison would not separate the two steps.

\paragraph{Setup.}
We use the 18 validator targets and the regexes of the main benchmark
(\S\ref{sec:benchmark}), with $N = 1000$ positives and $N = 1000$ negatives. Each
configuration runs 10 times with seeds $0$--$9$, and each run is paired with the
Full run that has the same target and seed. Shortest-first enumeration
is deterministic, so its positive set is the same in every run. For every run we
pass all positives and negatives to the target validator and measure line and
branch coverage of the validator's code with \texttt{coverage.py}, using the same
harness as in \S\ref{sec:rq-coverage}. For each target we report
$\Delta = \text{variant} - \text{Full}$ in percentage points, averaged over the 10
runs, so a negative $\Delta$ means that coverage is lost without the removed or
replaced component. Across targets we report the mean $\Delta$, a 95\% bootstrap confidence
interval over the 18 targets (10{,}000 resamples), and the number of targets where
Full is better, tied, or worse. A difference below $0.005$ points counts as a tie.
We bootstrap over targets rather than over the 180 target--run pairs, because runs
of the same target are strongly correlated.

\begin{table*}
\centering
\caption{Ablation on 18 validation targets (mean of 10 runs). $\Delta$ is the variant
minus Full in percentage points; negative means coverage is lost without the component. Only
targets with a non-zero $\Delta$ are listed; means are over all 18.}
\label{tab:ablation}
% \scriptsize
\setlength{\tabcolsep}{4pt}
\renewcommand{\arraystretch}{1.05}
\resizebox{\textwidth}{!}{%
\begin{tabular}{l rr rr rr rr rr}
\toprule
 & \multicolumn{2}{c}{Full} & \multicolumn{2}{c}{$\Delta$ Shortest-first enum.} & \multicolumn{2}{c}{$\Delta$ No positives} & \multicolumn{2}{c}{$\Delta$ Local complements only} & \multicolumn{2}{c}{$\Delta$ Mutations only} \\
\cmidrule(lr){2-3}\cmidrule(lr){4-5}\cmidrule(lr){6-7}\cmidrule(lr){8-9}\cmidrule(lr){10-11}
Target & Line & Branch & Line & Branch & Line & Branch & Line & Branch & Line & Branch \\
\midrule
datetime\_dateutil & 51.63 & 54.44 & \phantom{$-$}0.00 & \phantom{$-$}0.00 & \phantom{$-$}0.00 & \phantom{$-$}0.00 & $-$15.22 & $-$17.78 & \phantom{$-$}0.00 & \phantom{$-$}0.00 \\
datetime\_django\_utils & 21.82 & 27.78 & \phantom{$-$}0.00 & \phantom{$-$}0.00 & \phantom{$-$}0.00 & \phantom{$-$}0.00 & $-$1.82 & $-$5.56 & \phantom{$-$}0.00 & \phantom{$-$}0.00 \\
domain\_django\_core & 2.49 & 5.43 & \phantom{$-$}0.00 & $-$1.09 & \phantom{$-$}0.00 & $-$1.09 & $-$0.03 & $-$0.11 & \phantom{$-$}0.00 & \phantom{$-$}0.00 \\
domain\_stdlib\_idna & 9.63 & 13.21 & \phantom{$-$}0.00 & \phantom{$-$}0.00 & \phantom{$-$}0.00 & \phantom{$-$}0.00 & $-$1.38 & $-$1.89 & $-$0.92 & $-$0.94 \\
email\_django\_core & 3.74 & 9.78 & \phantom{$-$}0.00 & \phantom{$-$}0.00 & $-$0.31 & $-$2.17 & \phantom{$-$}0.00 & \phantom{$-$}0.00 & \phantom{$-$}0.00 & \phantom{$-$}0.00 \\
email\_validator\_package & 46.38 & 44.87 & $-$1.59 & $-$0.67 & $-$1.59 & $-$0.67 & $-$2.08 & $-$3.12 & $-$0.78 & $-$0.89 \\
ip\_stdlib\_ipaddress & 11.10 & 16.67 & \phantom{$-$}0.00 & \phantom{$-$}0.00 & \phantom{$-$}0.00 & \phantom{$-$}0.00 & $-$3.27 & $-$5.56 & $-$0.20 & $-$0.29 \\
url\_django\_core & 8.35 & 19.35 & $-$0.25 & $-$0.87 & $-$1.18 & $-$4.13 & $-$0.62 & $-$2.17 & \phantom{$-$}0.00 & \phantom{$-$}0.00 \\
url\_stdlib\_urllib & 10.63 & 13.60 & $-$0.18 & $-$0.43 & $-$0.18 & $-$0.43 & $-$1.29 & $-$0.78 & \phantom{$-$}0.00 & \phantom{$-$}0.00 \\
\midrule
Mean (18 targets) & 15.39 & 18.98 & $-$0.11 & $-$0.17 & $-$0.18 & $-$0.47 & $-$1.43 & $-$2.05 & $-$0.11 & $-$0.12 \\
\bottomrule
\end{tabular}%
}
\end{table*}


\paragraph{Results: boundary positives.}
Removing the positives lowers coverage by $0.18$ line and $0.47$ branch points on
average (95\% CI $[-1.01, -0.07]$ for branches). Full is better on 5 of the 18
targets and never worse. The losses are concentrated where the success path of the
validator contains real logic: \texttt{url\_django\_core} ($-4.13$ branch),
\texttt{email\_django\_core} ($-2.17$), \texttt{domain\_django\_core} ($-1.09$),
\texttt{email\_validator\_package} ($-0.67$), and \texttt{url\_stdlib\_urllib}
($-0.43$). Replacing the boundary positives with shortest-first enumeration recovers
part of this loss but not all of it ($\Delta = -0.11$ line, $-0.17$ branch; better on 4
targets, never worse). The remaining gap is largest on
\texttt{email\_validator\_package} ($-1.59$ line), where shortest-first enumeration reaches
no more coverage than having no positives at all. All 1000 of its strings have six
characters and the form \texttt{!@0.}$xy$ or \texttt{!@1.}$xy$, so it never reaches
multi-character local parts, dotted local parts, multi-label domains, long labels,
or the \texttt{xn--} branch. Boundary positives reach these by design.

The averages are small because, with the benchmark regexes, the negatives already
do much of the positives' work. Each benchmark regex is a practical approximation
of what its validator accepts, and most validators accept more than the regex
does (for example, \texttt{uuid.UUID} ignores hyphens and \texttt{ipaddress}
accepts compressed IPv6). Many strings that the regex rejects are therefore
accepted by the validator and exercise its success path. To isolate the
positives, we repeated this comparison with regexes written to match each
validator's actual acceptance rule (14 of the 18 targets; the other four keep the
benchmark regex). With these regexes, removing positives gives $\Delta = -3.02$ line and
$-4.62$ branch points, and replacing them with shortest-first enumeration gives
$-2.02$ and $-2.85$.

\paragraph{Results: local complements and context-aware mutations.}
The two negative steps contribute very differently. With local complements alone
(about 30 negatives per target), coverage falls by $1.43$ line and $2.05$ branch
points (95\% CI $[-4.35, -0.53]$), and Full is better on 8 of the 18 targets. The
largest losses are on parsers whose rejection logic branches on the shape of the
input: \texttt{datetime\_dateutil} ($-17.78$ branch), \texttt{ip\_stdlib\_ipaddress}
and \texttt{datetime\_django\_utils} ($-5.56$ each), \texttt{email\_validator\_package}
($-3.12$), and \texttt{url\_django\_core} ($-2.17$). A single wrong character
reaches the first rejection check of these validators, but not the checks behind
it, such as a wrong number of fields, an out-of-range component, or a misplaced
separator. Context-aware mutations reach these checks. On the other 10 targets,
step~1 alone already matches Full.

Without step~1's strings, context-aware mutations alone (about 202 negatives per
target) come within $0.11$ line and $0.12$ branch points of Full and are tied on
15 of the 18 targets. The remaining gaps are on \texttt{domain\_stdlib\_idna}
($-0.94$ branch), \texttt{email\_validator\_package} ($-0.89$), and
\texttt{ip\_stdlib\_ipaddress} ($-0.29$). These gaps cannot be attributed to step~1
alone, because Full also includes the fallback fill, which the step variants do
not. On \texttt{domain\_stdlib\_idna} the gap comes from strings that only the
fallback fill produces.

\paragraph{Why most $\Delta$ are zero.}
Coverage changes only when an input reaches a line that no other input reaches.
Most validators reject malformed input at one or two checks, so negatives with
very different faults end at the same line, and removing the component that
produced one of them changes nothing. Three of the targets with all-zero rows
show this (line numbers are for Python~3.13 and Django~6.0).
\begin{itemize}
  \item \texttt{uuid\_stdlib}. \texttt{uuid.UUID} removes hyphens, checks the
    length, and then parses the rest as hex. The local complement
    \texttt{00000000:0000-0000-0000-000000000000} (a colon instead of the first
    hyphen) and the whitespace mutation
    \texttt{\textvisiblespace 00000000-0000-0000-0000-000000000000} are different
    faults, but both leave 33 characters and stop at the same
    \texttt{raise ValueError('badly formed hexadecimal UUID string')}
    (\texttt{uuid.py}, line~181). A wrong character such as
    \texttt{!0000000-\ldots} fails one line later, in \texttt{int(hex, 16)}. These two
    lines catch all 37 local complements and 191 of the 203 mutations. The other
    12 mutations, such as \texttt{0000000-0000-0000-0000-0000000000000} with its
    hyphens in the wrong places, are accepted, because \texttt{uuid.UUID} ignores
    hyphens. They therefore follow the success path, which the positives already
    cover.
  \item \texttt{datetime\_stdlib\_strptime}. \texttt{strptime} compiles the whole
    format into one regex and reports every mismatch with the same error. The
    local complement \texttt{a900-01-01T00:00:00} (a letter for the first digit),
    the structure mutation \texttt{00:00:00T10-10-0091} (the string reversed), and
    the whitespace mutation \texttt{\textvisiblespace 1900-01-01T00:00:00} execute
    exactly the same lines and stop at
    \texttt{raise ValueError("time data \%r does not match format \%r")}
    (\texttt{\_strptime.py}, line~555). 26 of the 28 local complements and 201 of
    the 219 mutations end there.
  \item \texttt{ip\_django\_core}. \texttt{validate\_ipv46\_address} tries IPv4,
    then IPv6, and raises one error if both fail. The local complement
    \texttt{a.0.0.0}, the structure mutation \texttt{.0.0.00}, the numeric mutation
    \texttt{-1}, and the whitespace mutation \texttt{\textvisiblespace 0.0.0.0}
    execute identical lines and end at the same \texttt{raise ValidationError}
    (\texttt{validators.py}, line~331). All 59 local complements and 415 of the 421
    mutations take this path. The remaining six, compressed IPv6 addresses such as
    \texttt{0::0:0:0:0:0:0}, are accepted by Django and reach no line that Full's
    other inputs miss.
\end{itemize}
In each case the validator does not distinguish the faults, so once one
component has reached the rejecting line, the others add nothing that coverage
can measure.

Taken together, context-aware mutations reach almost all of the negative-side
coverage, and local complements add little once mutations are present. Local
complements are still useful as a cheap first step, because about 30 strings per
target reach Full's coverage on 10 of the 18 targets. The boundary positives
contribute most on validators whose success path has real logic, and the
contribution grows when the regex matches the validator closely.

\paragraph{Limitations.}
The two step variants use fewer negatives than Full (about 30 and 202 against 1000),
because we removed the fallback fill to keep the steps separate. Their $\Delta$
values therefore combine the effect of the missing step with the effect of a
smaller test set. Coverage is also a coarse measure here. As the examples above
show, negatives with different faults count the same once one of them has reached
the rejecting line, so an ablation based on coverage understates differences
between components that fail at the same check.
